\chapter{Methodology}

\section{Problem Definition and Label Taxonomy}
\section{Pipeline Overview}
% Include the end-to-end pipeline figure.
\section{Synthetic Data Generation}
\section{Solver--Critic Forward Pass}
\section{Interaction Labeling}
% Labeling iterations, Claude judge, prompt hash/settings, human audit, Cohen's
% kappa, and negative-class feedback audit.
\section{Activation Caching}
\section{Layer Selection Probe}
% Explain the probe mechanism, candidate layers, leakage-safe setup, metric
% choice, and why those metrics are appropriate. Report how the results led to
% the final layer rather than presenting layer choice as an assumption.
\section{SAE Training}

\subsection{Training Data and Preprocessing}

The custom SAEs were trained on the Solver's Layer 25 residual-post activations
from the second Solver turn. Each example is the 2,560-dimensional activation
at the final prompt token, immediately before the Solver generates its response
to the candidate answers. The training, validation, test, and intervention
partitions follow the question-family split defined before feature selection.
Only the training partition updates SAE parameters. Validation activations are
used for checkpoint and hyperparameter selection. The test partition is
evaluated only after model selection, and the intervention partition remains
outside SAE fitting and selection.

All inputs are normalized by one scalar root-mean-square value computed over
every element of the training activation matrix. The same frozen value is used
for all partitions, avoiding split-specific normalization. Decoder outputs are
rescaled to the original activation units before they are inserted into the
Solver. The full encoding, reconstruction, and training objectives are given
below.

\begin{align}
    s_{\mathrm{RMS}}
    &= \sqrt{\frac{1}{Nd}\sum_{i=1}^{N}\sum_{j=1}^{d}x_{ij}^{2}},
    & \tilde{x}_i &= \frac{x_i}{s_{\mathrm{RMS}}}, \\
    a_i
    &= W_{\mathrm{enc}}(\tilde{x}_i-b_{\mathrm{dec}})
       + b_{\mathrm{enc}},
    & z_i &= g(a_i), \\
    \hat{\tilde{x}}_i
    &= W_{\mathrm{dec}}z_i+b_{\mathrm{dec}},
    & \hat{x}_i &= s_{\mathrm{RMS}}\hat{\tilde{x}}_i .
\end{align}

For the $L_1$ model, $g(a)=\operatorname{ReLU}(a)$ and the training objective is

\begin{equation}
    \mathcal{L}_{L_1}
    = \frac{1}{Nd}\sum_{i=1}^{N}
      \left\lVert\tilde{x}_i-\hat{\tilde{x}}_i\right\rVert_2^2
      + \lambda\frac{1}{Nm}\sum_{i=1}^{N}\lVert z_i\rVert_1 .
\end{equation}

For the TopK model, $g(a)=\operatorname{TopK}(\operatorname{ReLU}(a),k)$ and

\begin{equation}
    \mathcal{L}_{\mathrm{TopK}}
    = \mathcal{L}_{\mathrm{rec}} + \alpha\mathcal{L}_{\mathrm{AuxK}},
\end{equation}

where $\mathcal{L}_{\mathrm{AuxK}}$ reconstructs the detached main-path residual
using eligible inactive features and is applied during training only.


Here, $N$ is the number of examples, $d=2{,}560$ is the residual-stream width,
$m=12{,}888$ is the custom dictionary width, and $z_i$ is the sparse latent
vector. Decoder columns are constrained to unit $L_2$ norm after each optimizer
step so that feature magnitude cannot be changed simply by increasing decoder
length.

\subsection{Custom SAE Training}

We trained two custom architectures with the same encoder--decoder dimensions
and data but different sparsity mechanisms. Both used tied encoder--decoder
initialization, Adam optimization, shuffled batches of 256, seed 42, and a
maximum of 200 epochs. A ReduceLROnPlateau scheduler reduced the learning rate
by a factor of ten after five validation epochs without improvement. Training
stopped after 20 epochs without an absolute validation-loss improvement greater
than $10^{-4}$, while the checkpoint with the lowest observed validation loss
was retained.

For the $L_1$ SAE, the encoder applies a nonnegative activation and the loss
adds the mean absolute latent activation with coefficient $\lambda$. The sweep
tested $\lambda\in\{1,2,3\}$ at learning rate $10^{-4}$; the selected model
uses $\lambda=2$. For the TopK SAE, the main path retains exactly the $k$
largest nonnegative activations per example. At learning rate
$3\times10^{-4}$, the sweep compared AuxK models with $k=64$ and $k=128$ and
a $k=128$ ablation without AuxK. The selected model uses $k=128$ and an AuxK
coefficient of $1/128$. Up to 512 features that had not activated
on the main path within the preceding 10,000 training examples were eligible
to reconstruct the detached residual. This auxiliary signal reduces the chance
that inactive features receive no useful gradient, but it does not contribute
to validation loss or activate features in the main TopK representation
\cite{gao2025scaling}.

We did not automatically select the lowest scalar loss across both sparsity
families because their losses are not directly comparable. Instead, one
checkpoint per family was chosen from its validation tradeoff among
reconstruction error, variance explained, mean active features, and dictionary
coverage. For a frozen checkpoint, variance explained and average activity are
defined as

\begin{equation}
    \operatorname{VE}
    = 1-
      \frac{\sum_i\lVert x_i-\hat{x}_i\rVert_2^2}
           {\sum_i\lVert x_i-\bar{x}_{\mathrm{train}}\rVert_2^2},
    \qquad
    L_0 = \frac{1}{N}\sum_{i=1}^{N}\lVert z_i\rVert_0 .
\end{equation}


where $\bar{x}_{\mathrm{train}}$ is the training-set mean activation. Feature
firing frequencies were also recorded as the fraction of examples for which a
latent was positive. After selection, the exact best checkpoint was reloaded
and evaluated without updating weights or inactivity counters.

TopK was retained as a training and representation-quality comparison but was
not carried into the final downstream feature-selection and causal tasking.
The final downstream configuration used the $L_1$ custom SAE and members of
the pretrained SAE family.

Algorithm~\ref{alg:custom-sae-training} summarizes the shared training loop and
the branch between the two sparsity mechanisms.

\begin{algorithm}[H]
\caption{Custom L1 or TopK--AuxK SAE training}
\label{alg:custom-sae-training}
\begin{algorithmic}[1]
\Require Training activations $X_{\mathrm{train}}$; validation activations $X_{\mathrm{val}}$;
sparsity mode $q\in\{L_1,\mathrm{TopK}\}$; dictionary width $m$; maximum epochs $E$;
batch size $B$; learning rate $\eta$
\Ensure Frozen SAE checkpoint with the lowest main-path validation loss
\State $s_{\mathrm{RMS}}\gets\sqrt{\operatorname{mean}(X_{\mathrm{train}}^2)}$
\State Initialize tied encoder--decoder weights and normalize every decoder column
\State Initialize Adam, the plateau scheduler, inactivity counters, and early stopping
\For{$e=1$ to $E$}
    \For{each shuffled minibatch $X\subset X_{\mathrm{train}}$ of size $B$}
        \State $\widetilde{X}\gets X/s_{\mathrm{RMS}}$
        \State $A\gets W_{\mathrm{enc}}\widetilde{X}+b_{\mathrm{enc}}$
        \If{$q=L_1$}
            \State $Z\gets\operatorname{ReLU}(A)$
            \State $\mathcal{L}\gets\mathcal{L}_{\mathrm{rec}}+\lambda\operatorname{mean}(|Z|)$
        \Else
            \State $P\gets\operatorname{ReLU}(A)$; $Z\gets\operatorname{TopK}(P,k)$
            \State Update each feature's last main-path firing position from $Z$
            \State $J\gets$ features inactive for at least $w$ training examples
            \If{$J\ne\emptyset$ and the inactivity warm-up is complete}
                \State $Z_{\mathrm{aux}}\gets
                \operatorname{TopK}(P|_J,\min(k_{\mathrm{aux}},|J|))$
                \State $R\gets(\widetilde{X}-\widehat{\widetilde{X}})\text{.detach}()$
                \State $\mathcal{L}\gets\mathcal{L}_{\mathrm{rec}}+
                \alpha\operatorname{MSE}(W_{\mathrm{dec}}Z_{\mathrm{aux}},R)$
            \Else
                \State $\mathcal{L}\gets\mathcal{L}_{\mathrm{rec}}$
            \EndIf
        \EndIf
        \State Update SAE parameters with Adam using $\nabla\mathcal{L}$
        \State Renormalize every decoder column to unit $L_2$ norm
    \EndFor
    \State Evaluate the main-path objective on $X_{\mathrm{val}}$ without AuxK
    \State Update the scheduler and save any new lowest-loss checkpoint
    \If{validation loss has no significant improvement for the patience window}
        \State \textbf{break}
    \EndIf
\EndFor
\State Reload and freeze the lowest-loss checkpoint
\State \Return frozen SAE and split-level reconstruction and activity metrics
\end{algorithmic}
\end{algorithm}


\subsection{Pretrained SAE and Decoder-Only Adaptation}

The external baseline is the frozen Layer 25 Gemma Scope 2 JumpReLU SAE
described in Chapter~4. Because it was trained on a broader distribution, we
also ran a decoder-only adaptation experiment. The pretrained encoder weights,
encoder bias, and JumpReLU thresholds were frozen; their latent activations were
cached once for the training and validation partitions. Only the decoder and
reconstruction bias were optimized against the interaction activations. This
preserves feature firing patterns and probe inputs while testing whether a
decoder fitted to the study distribution reconstructs and steers more
effectively.

Decoder adaptation used seed 42, batches of 256, learning rate $10^{-5}$, a
maximum of 10 epochs, and early-stopping patience of five. The run manifest
records SHA-256 hashes of the source SAE parameters, activation tensors, probe
configuration, and cached encoder features. A separate full-fine-tuning mode
was implemented as an exploratory comparison, but updating the encoder changes
the feature basis and therefore weakens one-to-one comparison with the
pretrained feature identities. The decoder-only model is the cleaner adaptation
for analyses that require the same selected latents before and after fitting.

Algorithm~\ref{alg:decoder-only-adaptation} gives the adaptation and integrity
checks used for the pretrained SAE.

\begin{algorithm}[H]
\caption{Decoder-only adaptation of the pretrained SAE}
\label{alg:decoder-only-adaptation}
\begin{algorithmic}[1]
\Require Pretrained JumpReLU SAE $f$; training activations $X_{\mathrm{train}}$;
validation activations $X_{\mathrm{val}}$; epochs $E$; batch size $B$; learning rate $\eta$
\Ensure Adapted decoder with unchanged encoder features and thresholds
\State Load $f$ and digest $(W_{\mathrm{enc}},b_{\mathrm{enc}},\theta)$
\State Freeze $W_{\mathrm{enc}}$, $b_{\mathrm{enc}}$, and JumpReLU thresholds $\theta$
\For{$S\in\{X_{\mathrm{train}},X_{\mathrm{val}}\}$}
    \State Cache $Z_S\gets\operatorname{JumpReLU}(W_{\mathrm{enc}}S+b_{\mathrm{enc}};\theta)$ once
\EndFor
\State Evaluate the untouched SAE and bias-only recentered baseline
\State Initialize Adam over only $(W_{\mathrm{dec}},b_{\mathrm{dec}})$
\State Save the untouched model as eligible epoch 0
\For{$e=1$ to $E$}
    \For{each shuffled minibatch $(Z,X)$ from the cached training codes}
        \State $\widehat{X}\gets W_{\mathrm{dec}}Z+b_{\mathrm{dec}}$
        \State $\mathcal{L}\gets\operatorname{MSE}(\widehat{X},X)$
        \State Update only $(W_{\mathrm{dec}},b_{\mathrm{dec}})$ with Adam
    \EndFor
    \State Evaluate validation reconstruction using the cached validation codes
    \If{validation MSE is lower than the best eligible MSE}
        \State Save the decoder and reconstruction bias; reset patience
    \Else
        \State Increment the early-stopping counter
    \EndIf
    \If{the counter reaches the patience limit}
        \State \textbf{break}
    \EndIf
\EndFor
\State Reload the best checkpoint and verify the encoder digest is unchanged
\State \Return adapted SAE, provenance manifest, and before--after metrics
\end{algorithmic}
\end{algorithm}


\subsection{Exploratory LLM-Guided SAE Editing}

Inspired by the MIT CSAIL ``Sparks of Weight-Space Programming'' demonstration
\cite{takehi2026sparks}, we conducted a post-hoc side experiment asking whether
an LLM could improve causal steering by directly editing selected SAE decoder
directions without gradient training. This experiment is exploratory and is
not included in the prespecified primary comparison.

Claude received train-partition measurements and a locked-down Python sandbox,
but no validation, test, intervention, or final causal results. It could inspect
the pretrained SAE and the 10 Adoption and 10 Retention features selected by
the probe, then propose edits to only those 20 decoder rows. Encoder parameters,
thresholds, decoder-row lengths, and every unselected row remained fixed. Each
edited row was constrained to cosine similarity of at least 0.95 with its
original direction. No optimizer, backpropagation, or automated parameter
search was permitted.

Two controls separated LLM reasoning from generic weight perturbation. A
mechanical edit rotated the same rows toward class-mean activation directions
using a fixed formula, and a random edit applied the same row-wise rotation
angles toward random orthogonal directions. All three edited SAEs then entered
the same held-out feature-selection and intervention pipeline. This design
tests whether an LLM-designed direction offers value beyond a deterministic
behavioral direction or an equal-size random modification; it does not imply
that the resulting decoder features acquired a complete semantic
interpretation.

The constrained interaction, authoritative host checks, and construction of
the two matched controls are summarized in
Algorithm~\ref{alg:llm-guided-sae-editing}.

\begin{algorithm}[H]
\caption{Constrained LLM-guided SAE decoder editing}
\label{alg:llm-guided-sae-editing}
\begin{algorithmic}[1]
\Require Original decoder $W_{\mathrm{dec}}^{(0)}$; 10 Adoption features $A$;
10 Retention features $R$; train-only sandbox data; action and steering budgets
\Ensure Auditable edited SAE and matched mechanical and random controls
\State Lock the encoder, thresholds, unselected rows, and decoder-row lengths
\State Start a no-network sandbox with train data, fixed codes, diagnostics, and \textsc{ApplyEdit}
\State Give the LLM the task and constraints, but no validation, test, intervention, or control results
\While{the LLM has not returned \textsc{Done} and no budget limit is reached}
    \State Receive one action: a Python snippet, \textsc{SteerTest}, or \textsc{Done}
    \If{the action proposes a row edit $(j,w_j')$}
        \State Check $j\in A\cup R$, finite values, and the required row shape
        \State Check $|\lVert w_j'\rVert_2/\lVert w_j^{(0)}\rVert_2-1|\le10^{-4}$
        \State Check $\cos(w_j',w_j^{(0)})\ge0.95$
        \If{every constraint passes}
            \State Accept the edit; record its rationale and hash; report diagnostics
        \Else
            \State Reject the edit and leave the authoritative decoder unchanged
        \EndIf
    \ElsIf{the action is \textsc{SteerTest} and budget remains}
        \State Test the current edit on the fixed train-only steering cohort
        \State Return flip and candidate-margin summaries to the LLM
    \Else
        \State Execute the permitted diagnostic snippet or record a protocol error
    \EndIf
\EndWhile
\State Replay accepted edits onto $W_{\mathrm{dec}}^{(0)}$ and save an auditable SAE
\State Build a mechanical control by rotating the same rows toward the class-mean direction
\State Build a random control using the LLM edit's angles and random orthogonal directions
\State Verify that both controls satisfy the same length and cosine constraints
\State \Return LLM-edited, mechanical-control, and random-control SAEs
\end{algorithmic}
\end{algorithm}

\section{Feature Selection Probe}
% Probe, ablation check, selectivity controls, random draws, and dominant feature.
\section{Causality Experiments}
% Initial zeroing/ablation; suppression and amplification; edit-norm controls;
% mechanical versus LLM-guided edits; raw-activation controls; and
% reconstruction control.
\section{Metrics and Statistical Analysis}
% Flip definition, confidence intervals, random baselines,
% and multiple-comparison handling.
\section{Experimental Setup and Reproducibility}

\subsection{Hardware and Software Environment}

Individual GPU jobs for collection, activation extraction, SAE training, and
causal replay used NVIDIA A10G devices with 23~GiB of device memory. Collection
and causal intervention were distributed across multiple workers; each worker
used one GPU. The
recorded machine used four CPU cores and approximately 15.4~GiB of system
memory for the exploratory SAE-editing workflow. The project environment is
defined for Python~3.11 and PyTorch~2.13, with Transformers~5.x, Datasets,
scikit-learn, SciPy, pandas, NumPy, and SAE Lens. Rather than relying only on
the environment specification, individual run manifests record the resolved
package versions; for example, the final collection provenance records
PyTorch~2.13.0, Transformers~5.17.0, and Datasets~5.0.1.

Model inference used bfloat16 on the GPU. The model and dataset sources were
pinned to immutable Hugging Face revisions:

\begin{itemize}
    \item Gemma Solver: \texttt{093f9f388b31de276ce2de164bdc2081324b9767};
    \item Qwen Critic: \texttt{cdbee75f17c01a7cc42f958dc650907174af0554};
    \item MuSiQue: \texttt{c8f4f8c9465fb69d31a8eae894c3fd509c4ca321}.
\end{itemize}

Causal replay used greedy decoding with at most 64 new tokens, ensuring that a
fixed prompt and model state could be checked against the originally stored
response.

\subsection{Randomness and Partition Integrity}

Seed 42 was fixed for collection sampling, SAE initialization and shuffling,
probe resampling, random-feature controls, and intervention ordering unless a
run explicitly records a different seed. The final dataset split used seed 43.
Question families, rather than
individual episodes, define the train, validation, test, and intervention
partitions. The code treats test and intervention partitions as sealed during
layer and feature selection. Edit-size selection used a designated 50-episode
pilot at intervention positions 1,001--1,050; those episodes were disjoint from
the 1,000-episode intervention study cohort. The study cohort was therefore not
used to choose a layer, SAE checkpoint, feature set, edit direction, or edit
size.

Random controls are stored or deterministically regenerated from the frozen
seed and episode identifiers. Completed outputs are resumed rather than
regenerated, which prevents partial distributed runs from silently changing
previous model responses. Evaluation scripts also record excluded episodes and
technical labeling failures instead of replacing them with newly sampled
examples.

\subsection{Configuration and Provenance Records}

Every major stage is driven by a versioned YAML configuration. Output
directories are run-specific and normally refuse to overwrite an existing
run. Run records store the complete command, Git commit, hashes of uncommitted
tracked code, configuration hash, model and SAE specifications, label-policy
version, and SHA-256 hashes of the input artifacts. Feature-selection outputs
also hash every SAE file they read. Consequently, a later causal run can verify
that the selected features refer to the same checkpoint rather than another
file with the same human-readable name.

Collection provenance records the ordered question manifest, dataset revision,
resolved model revisions, package versions, and available GPUs. SAE runs store
their full hyperparameters, learning-rate history, best epoch, checkpoint,
split-level reconstruction metrics, and feature firing frequencies. Causal
runs additionally freeze the episode IDs, conditions, feature sets, random
draws, edit size, and analysis settings before distributed workers begin.

\subsection{Run-Consistency Checks}

The causal pipeline performs a plain replay before applying any activation
edit. It compares the regenerated prompt hash, response, first token, and
behavior label with the stored collection artifacts. If baseline agreement
falls below the configured threshold, the run stops before generating edited
conditions. Distributed workers must reproduce the successful preflight; shard
sets with missing, duplicated, or incompatible conditions are rejected during
merge.

The pipeline also rejects mismatched SAE hashes, changed label protocols,
incompatible feature-reference files, and arms generated under different
frozen settings. Generated rows record the actual edit norm, reconstruction
error, number of active and edited features, token budget status, and elapsed
generation time. These checks do not remove all sources of nondeterminism, but
they make discrepancies visible and prevent results from different model,
code, label, or intervention states from being combined silently.

The repository contains the configurations, scripts, lightweight summaries,
and exact run identifiers needed to trace every reported result. Large model
weights, activation tensors, and SAE checkpoints are stored separately because
of their size; reproduction requires retrieving the artifacts whose hashes are
recorded in the manifests. Appendix~E provides the command sequence and
artifact locations used to reconstruct the analysis.
